% Generated draft: requires mathematical and source review.
\documentclass[11pt]{article}
\usepackage[margin=25mm]{geometry}
\usepackage{fontspec}
\setmainfont{Latin Modern Roman}
\usepackage{amsmath,amssymb,mathtools,unicode-math}
\setmathfont{Latin Modern Math}
\usepackage{enumitem,needspace,xurl,hyperref}
\setlength{\parindent}{0pt}
\setlength{\parskip}{5pt}
\setlength{\emergencystretch}{3em}
\newcommand{\N}{\mathbb{N}}
\newcommand{\Z}{\mathbb{Z}}
\newcommand{\Q}{\mathbb{Q}}
\newcommand{\R}{\mathbb{R}}
\newcommand{\C}{\mathbb{C}}
\newcommand{\F}{\mathbb{F}}
\newcommand{\A}{\mathbb{A}}
\newcommand{\PP}{\mathbb{P}}
\newcommand{\E}{\mathbb{E}}
\newcommand{\eps}{\varepsilon}
\newcommand{\dd}{\,\mathrm{d}}
\begin{document}
\section*{UnsolvedMath ID: 11100028}
\textbf{Source problem code:} catalogue entry
\subsection{Definitions and mathematical statement}
In the standard closed smooth setting, let \(M\) be a spin manifold of dimension \(d\geq5\). The spin orientation \(M\operatorname{Spin}\to KO\) gives \(\alpha(M)\in KO_d(\mathrm{pt})\). The source reports Stefan Stolz's theorem that a simply connected such \(M\) admits a metric of everywhere positive scalar curvature exactly when \(\alpha(M)=0\).

Let \(B\) be a Bott manifold: in the usual convention it is a closed spin \(8\)-manifold whose index is the periodic Bott generator \(\beta\in KO_8(\mathrm{pt})\). The final printed conjecture asks, without an index hypothesis, whether every spin \(M\) of dimension at least \(5\) has some finite \(k\geq0\) for which \(M\times B^k\) admits a positive-scalar-curvature metric.

That missing hypothesis is material. Index multiplicativity gives
\[
 \alpha(M\times B^k)=\alpha(M)\beta^k,
\]
and Bott multiplication is invertible in periodic \(KO\), while positive scalar curvature forces this index to vanish. Thus the unconditional printed version cannot bypass the index obstruction. The summary instead describes a conditional stable problem with vanishing \(\alpha\)-invariant; this is retained as a possible intended correction, not silently inserted.

For nontrivial fundamental group \(\pi\), the relevant obstruction may be a higher Rosenberg index in real group-\(C^*\)-algebra \(KO\), rather than merely \(\alpha(M)\) at a point. The summary does not specify reduced/maximal algebra or exact group/assembly hypotheses. It reports stable equivalence under a strong-Novikov-type assumption and leaves general groups open, while the source recalls counterexamples to an unstabilized formulation. Those are contextual reports without precise citations, not supplied unrestricted theorems.
\subsection{Short English statement}
Study positive scalar curvature after multiplying by Bott manifolds, while explicitly recording that the printed universal claim omits an index-vanishing condition and that nontrivial fundamental groups require the correct higher index.
\subsection{Sources}
\begin{itemize}
\item Mark Hovey, \emph{Algebraic Topology Problem List}, ``Applications,'' item 4, surviving mirror \url{https://www-users.cse.umn.edu/~tlawson/hovey/}. Recorded access 2026-07-29; extraction html-ordered-list.
\item Stefan Stolz, source attribution for the simply connected spin theorem; the Novikov conjecture and unstabilized counterexamples are mentioned without complete references. The summary gives no exact reference for its strong-Novikov conditional stable result.
\item Local UnsolvedMath registry, \texttt{data/unsolvedmath/problems.json}, permanent ID \(11100028\), exact source code \texttt{catalogue entry}; supplied partially\_solved/PARTIAL-PROGRESS status.
\end{itemize}
\end{document}
